\section{第2组：质点动力学}
\sourceproblem{2-T1}{3}{7}
\problemtext 静止释放质量 $m$ 的物体，受线性阻力 $-kv$，求下落速度与终速。
\workings 取向下为正，$m\dot v=mg-kv$，$v(0)=0$。积分得 $v=mg(1-e^{-kt/m})/k$，终速 $mg/k$。
\crosscheck $k>0$；$t=0$ 时 $\dot v=g$，$k\to0$ 时 $v\to gt$。
\sourceproblem{2-T2}{3}{7}
\problemtext $m=2.00\,\mathrm{kg}$，$F=10+6x^2$（SI），从 $x=0$ 静止出发，求 $x=4.0\,\mathrm m$ 速率。
\workings 牛顿定律 $mv\,\mathrm dv/\mathrm dx=10+6x^2$；积分 $mv^2/2=10x+2x^3$，故 $v=\sqrt{168}\,\mathrm{m/s}\approx12.96\,\mathrm{m/s}$。
\crosscheck 功能关系独立核验：外力做功 $168\,\mathrm J$，等于动能增量；保留正根。
\sourceproblem{2-T3}{3}{7}
\problemtext 质量 $m$ 以 $v_0>0$ 竖直上抛，阻力与速度线性相关，求速度及最大上升高度。
\workings 向上为正，$m\dot v=-mg-kv$，解得
\[
v=(v_0+mg/k)e^{-kt/m}-mg/k.
\]
令 $v=0$ 得 $t_*=(m/k)\ln(1+kv_0/(mg))$。积分至该时刻，
\[
H=\frac{mv_0}{k}-\frac{m^2g}{k^2}\ln(1+kv_0/(mg)).
\]
\crosscheck 用 $\ln(1+z)=z-z^2/2+O(z^3)$ 检验 $k\to0$ 时 $H\to v_0^2/(2g)$；高度为相对抛出点。
\sourceproblem{2-T4}{4}{8}
\problemtext 快艇关机时速度 $v_0$，阻力大小 $kv^2$，质量 $m$，求速度、路程随时间及速度随路程的规律。
\workings 沿初速度取正，$m\dot v=-kv^2$，得
\[
v=\frac{v_0}{1+kv_0t/m},\quad
s=\frac mk\ln(1+kv_0t/m),\quad v=v_0e^{-ks/m}.
\]
\crosscheck $s(0)=0$，$\dot s=v$；速度始终为正并趋零，纯平方阻力模型下无有限停车时间。
\sourceproblem{2-T5}{4}{8}
\problemtext 水平加速度为 $a$ 的车厢内，悬挂小球相对车厢静止，求绳偏离竖直的角。
\workings 车厢系引入反向惯性力 $ma$，绳张力满足 $T\sin\theta=ma,T\cos\theta=mg$，故 $\theta=\arctan(a/g)$，绳向车加速度的反方向偏斜。
\crosscheck 地面惯性系同样有水平合力 $ma$、竖直合力零；质量约去。
\sourceproblem{2-T6}{4}{8}
\problemtext 上升加速度 $a$ 的电梯内，理想滑轮两侧悬挂 $m_1>m_2$，求张力及 $m_1$ 相对电梯的加速度。
\workings 在电梯系取有效重力 $g'=g+a$。令 $m_1$ 向下相对加速度 $b$，则 $m_1g'-T=m_1b$，$T-m_2g'=m_2b$，联立得
\[
b=\frac{m_1-m_2}{m_1+m_2}(g+a),\qquad
T=\frac{2m_1m_2}{m_1+m_2}(g+a).
\]
\crosscheck $a=0$ 退化为标准 Atwood 机，$m_1=m_2$ 时相对加速度零；假定绳绷紧。
\sourceproblem{2-T7}{5}{9}
\problemtext 教材一维冲量例题：力随时间按 $F=600-2\times10^5t$（SI）线性减小，末态速度 $300\,\mathrm{m/s}$，末态力为零；求作用时长、冲量与质量。
\workings $F(t_*)=0$ 得 $t_*=0.003\,\mathrm s$；冲量为三角形面积 $I=600t_*/2=0.9\,\mathrm{N\,s}$。在初速度为零的通常约定下，$m=I/300=0.003\,\mathrm{kg}$。
\crosscheck 若初速度未约定为零，则只能得 $m(v_f-v_i)=I$；质量数值依赖该初始条件。
\sourceproblem{2-T8}{5}{9}
\problemtext 稳定水流经直角弯管，质量流率 $3\times10^3\,\mathrm{kg/s}$、速率 $10\,\mathrm{m/s}$，求水流改变方向产生的反作用力。
\workings 入口、出口单位方向分别为 $\bm e_i,\bm e_o$，互相垂直。动量流变化率为 $\dot m v(\bm e_o-\bm e_i)$，其反向为 $\dot m v(\bm e_i-\bm e_o)$，模 $3\sqrt2\times10^4\,\mathrm N$，沿入口速度与反向出口速度的角平分线。
\crosscheck 这是忽略端面压力和流体重量的动量流反作用结果；若求真实管壁总受力，需另给端面压力、截面积与重力等条件，不能仅凭流量确定。
\sourceproblem{2-T9}{5}{9}
\problemtext 煤粉从高 $h=0.5\,\mathrm m$ 自由落至向右速率 $u=2.0\,\mathrm{m/s}$ 的带面；质量流率 $q_m=40\,\mathrm{kg/s}$，忽略已随带运动的积煤重量，求煤对带作用力。
\workings 取右和上为正，入口速度 $(0,-\sqrt{2gh})$，出口速度 $(u,0)$。忽略接触区滞留重量，带对煤的动量流作用为 $q_m(u,\sqrt{2gh})$，故煤对带力
\[
\bm F=-q_mu\,\bm i-q_m\sqrt{2gh}\,\bm j.
\]
取 $g=9.8$，得向左 $80\,\mathrm N$、向下约 $125.2\,\mathrm N$，模约 $148.6\,\mathrm N$，左下方与水平夹角 $\arctan(\sqrt{2gh}/u)\approx57.4^\circ$。
\crosscheck 水平力阻碍带运动，竖直力向下；若计入积煤或滞留煤重量，竖直负载另加。

\sourceproblem{2-T10}{6}{10}
\problemtext 均匀柔绳线密度 $\lambda$，从水平面以恒定速率 $v$ 竖直提起，求提起长度 $L$ 时拉力。
\workings 取正在上升的绳段为开系统，竖直动量 $\lambda Lv$，其增长率 $\lambda v^2$；进入绳段前竖直速度为零。故
\[
F-\lambda Lg=\lambda v^2,\qquad F=\lambda(Lg+v^2).
\]
\crosscheck 重力项随长度增加，动量通量项恒定；假定绳在提起处从静止转入竖直运动、无额外支持力作用于运动段。不可把变质量系统直接当固定质量系统。
\sourceproblem{2-T11}{6}{10}
\problemtext 卫星近地点高度439千米、远地点高度2384千米，地球半径6378千米，近地点速率8.10千米每秒，求远地点速率。
\workings 地心引力为中心力，近远地点速度垂直径矢；故
\[
v_2=v_1\frac{R+h_1}{R+h_2}
=8.10\frac{6817}{8762}\,\mathrm{km/s}\approx6.30\,\mathrm{km/s}.
\]
\crosscheck 应使用地心距离而非离地高度；远地点速率较小。
\sourceproblem{2-T12}{7}{11}
\problemtext 判断保守力正功是否增加势能、闭路保守力功、作用反作用力总功。
\workings $W_c=-\Delta U$，故正功使势能减少；闭路功为零；第三项因两物体位移可不同而不成立。仅第二项正确，选C。
\crosscheck 与天津理工第2章选择1的命题第一项措辞不同，但结论相同；两书来源编号分别保留。
\sourceproblem{2-T13}{7}{11}
\problemtext 恒力 $\bm F=4\bm i+3\bm j$（N），从原点至 $(8,6)$ 米，分别沿直线、先横后竖、先竖后横路径求功。
\workings 直线：$W=4(8)+3(6)=50\,\mathrm J$。先横后竖两段分别32和18焦耳，总50焦耳；先竖后横分别18和32焦耳，总50焦耳。任意路径的积分为 $\int(4\,\mathrm dx+3\,\mathrm dy)=4\Delta x+3\Delta y$，故为保守力，势能可取 $U=-4x-3y+C$（SI）。
\crosscheck 不仅三条指定路径相同，而且对任意可微路径成立，故保守性得到证明。
\sourceproblem{2-T14}{7}{11}
\problemtext $\bm r=a\cos\omega t\,\bm i+b\sin\omega t\,\bm j$，质量 $m$，$a>b>0$，求 $A(a,0)$、$B(0,b)$ 处动能、受力、从A到B的力及分力功，并判断保守性。
\workings $\bm v=\omega(-a\sin\omega t\,\bm i+b\cos\omega t\,\bm j)$，故 $K_A=mb^2\omega^2/2,K_B=ma^2\omega^2/2$。$\bm F=-m\omega^2(x\bm i+y\bm j)$；
\[
W_x=-m\omega^2\int_a^0 x\,\mathrm dx=\tfrac12m\omega^2a^2,\qquad
W_y=-m\omega^2\int_0^b y\,\mathrm dy=-\tfrac12m\omega^2b^2.
\]
总功 $m\omega^2(a^2-b^2)/2$。若将此力推广为整个平面上只依赖位置的上述中心力，则势能为 $m\omega^2(x^2+y^2)/2$，为保守力。
\crosscheck 总功等于动能差；仅凭单条轨迹不能唯一确定全空间力场，保守性判断明确采用位置力场模型。
\sourceproblem{2-T15}{8}{12}
\problemtext 光滑半球碗半径 $r$，质点从与向下竖直轴成 $\theta_0$ 的位置沿内面切线水平发射，求恰达水平碗口的初速。
\workings 竖直轴角动量守恒 $mr\sin\theta_0\,v_0=mr v_{\rm rim}$，临界到达碗口时沿经线速度为零。机械能守恒
\[
\tfrac12mv_0^2=\tfrac12mv_0^2\sin^2\theta_0+mgr\cos\theta_0.
\]
故 $v_0=\sqrt{2gr/\cos\theta_0}$，适用 $0\le\theta_0<\pi/2$。为验证途中可达，令 $\theta\ge\theta_0$，经线速度平方为
\[
v_\theta^2=v_0^2(1-\sin^2\theta_0/\sin^2\theta)-2gr(\cos\theta_0-\cos\theta).
\]
临界值代入后，令 $u=\cos\theta,u_0=\cos\theta_0$，对 $0<u<u_0$ 有
\[
v_\theta^2=2gr\,u\,\frac{(u_0-u)(1+u_0u)}{u_0(1-u^2)} >0.
\]
故中间没有势垒阻挡；底点情形由经线运动直接得到。
\crosscheck 光滑内面约束力不做功且对竖直轴力矩为零；碗口极限需单独处理，不能直接令 $\cos\theta_0=0$。
\sourceproblem{2-T16}{8}{12}
\problemtext 飞船在半径 $R_0$ 的圆轨道上速率为 $v_0$，变轨至近心距 $R_0$、远心距 $3R_0$ 的椭圆，求变轨后的近心点速率。
\workings 设中心天体引力参数 $\mu$，圆轨道 $v_0^2=\mu/R_0$。椭圆半长轴 $a=2R_0$，由活力公式
\[
v_p^2=\mu(2/R_0-1/(2R_0))=\tfrac32v_0^2,\qquad v_p=\sqrt{3/2}\,v_0.
\]
\crosscheck 可用近远心点角动量关系 $v_a=v_p/3$ 与机械能守恒独立推出同一结果；近心点变轨需加速。
